\chapter{Static solution methods}\label{ch:static}

\section{Which approximation answers which question?}\label{sec:method-landscape}

Representing the reservoir continuum and solving the interacting many-body
problem require separate approximations and checks. Exact diagonalization
solves a specified finite bath, whose continuum accuracy is assessed by bath
refinement. Quadratic continuum methods describe hybridization and pairing;
an interacting dot requires an additional treatment of its correlations.

\begin{table}[htb]
\centering\small
\caption{A guide to static methods. Each row describes a physical approximation
or representation. The QP and DMRG implementations are described below.}
\label{tab:method-landscape}
\begin{tabularx}{\textwidth}{p{.20\textwidth}XX}
\toprule
Method & Useful physical information & Validity and convergence \\
\midrule
Scattering / BdG & ABS, electron--hole coherence, phase-dependent current & Quadratic model; the short-contact formula also assumes short dwell time \\
Large-gap dot & Pair mixing, parity boundary, analytic matrix elements & Infinite-gap limit with lead quasiparticles integrated out \\
Perturbation / fRG & Virtual processes, weak-coupling corrections & Specify the expansion or truncation and its validity \\
NRG & Kondo scales, impurity thermodynamics, subgap spectra & Logarithmic bath and retained-state convergence \\
QMC & Equilibrium correlations and finite-temperature observables & Sampling errors; real-frequency continuation is difficult \\
ED / QP expansion & Many-body states, energies and transitions & Finite bath; QP truncation is an additional approximation \\
DMRG / MPS & Correlated states with entanglement-controlled compression & Bond dimension, optimization and bath convergence \\
\bottomrule
\end{tabularx}
\end{table}

\subsection{Scattering theory, BdG and effective dot models}

For a short, non-interacting contact, scattering theory reduces the normal
region to its transmission channels. The result \cref{eq:contact-abs} supplies
a direct analytic benchmark. A more general quadratic BdG Hamiltonian can
include spatial structure, multiple orbitals, Zeeman terms and spin--orbit
coupling. Diagonalizing that specified quadratic problem is exact. Replacing
the interacting dot term by a static mean field is a separate approximation,
which can miss quantum spin correlations or produce broken-symmetry solutions
where the finite impurity has a correlated singlet. Bulk BCS mean field describes
the reservoir condensate; mean-field decoupling of the dot's $U$ approximates its local correlations.

The large-gap model is another controlled simplification with a different
purpose: it retains the dot interaction exactly while integrating out the lead
quasiparticles. Its small Hilbert space exposes pair mixing and the parity
transition. It is complementary to a short resonant-level junction. A broad
resonance with $\Gamma\gg\Delta$ can satisfy the short-dwell-time condition,
whereas the conventional large-gap dot assumes $\Delta$ exceeds the dot and
hybridization scales after the wide-band limit. Choose between these approximations
according to the hierarchy of scales at the parameter point of interest.

\subsection{Perturbative and renormalization approaches}

An expansion in tunnelling describes virtual charge transfer and Josephson
cotunnelling in a weakly coupled dot. An expansion in $U$ starts instead from
a hybridized non-interacting problem. The expansion parameters and regimes of validity differ.
Near a charge degeneracy, denominators in a tunnelling expansion become small;
at low normal-state energies, Kondo logarithms can invalidate a finite-order
exchange expansion even when its bare coupling is small. Use a few virtual processes
to interpret the mechanism, and establish their quantitative accuracy by checking
the expansion in the relevant parameter regime.

Functional renormalization group (fRG) follows the evolution of propagators and
interaction vertices as a cutoff is removed. Practical truncations can describe
substantial parameter ranges efficiently, but the retained frequency dependence
and vertex hierarchy determine their range of validity. The dot-junction reviews
\cite{martinrodero2011josephson,meden2019andersonjosephson} compare these methods
and their limitations. Assess both perturbative and nonperturbative calculations
through their assumptions and convergence at the parameter point of interest.

\subsection{Numerical renormalization group}

NRG addresses an impurity coupled to a continuum with widely separated energy
scales~\cite{wilson1975renormalization,bulla2008numerical}. A logarithmic
discretization is transformed to a chain with successively smaller hopping
scales. The chain is diagonalized iteratively, retaining low-energy many-body
states at each step. The scale separation is what permits the buildup of Kondo
correlations to be followed without diagonalizing the entire Fock space at once.
A superconducting gap modifies this low-energy structure and produces discrete
subgap levels.

NRG is an important reference for the superconducting Anderson impurity,
including singlet--doublet boundaries and quantum YSR spectra. Its logarithmic
discretization, retained-state number and spectral processing still require
checks. More orbitals or channels increase the computational cost rapidly.
Broadening a discrete numerical spectrum is useful for visualization; predicting
physical lifetimes requires a dynamical calculation. For the isolated subgap poles, energies and
weights are often the most direct comparison with other methods.

\subsection{Quantum Monte Carlo}

Quantum Monte Carlo (QMC) samples an interacting equilibrium problem statistically,
often through imaginary-time correlation functions. Continuous-time impurity
algorithms can sum an expansion to all orders by stochastic sampling. In
``weak-coupling CT-QMC'', the name identifies the interaction expansion whose
orders are sampled~\cite{gull2011continuous}.

QMC provides a complementary route to finite-temperature occupation, current and
correlations. Statistical uncertainty, low-temperature cost and possible
model-dependent sign or phase problems must be assessed. Recovering a sharp
real-frequency spectrum from noisy imaginary-time data is an ill-conditioned
analytic-continuation problem. A method can therefore determine a thermal
current accurately while resolving nearby subgap spectral lines less directly.

\begin{figure}[htb]
\centering\manualfigure{primer_methods}
\caption{Reservoir representations and many-body approximations. A finite star
couples bath orbitals directly to the dot; a unitary transformation can organize
the same finite one-body bath as a chain. This coordinate change is exact.
NRG combines a logarithmic discretization with iterative
truncation; ED, QP truncation and MPS select different representations of the
many-body state. The orbital chain is a computational representation of the
reservoir's contact spectral measure.}
\label{fig:primer-methods}
\end{figure}

\subsection{Choosing a comparison}

For a quadratic junction, compare with an independent electron-space BdG
calculation. For the large-gap dot, compare with its exact even and odd blocks.
For a small finite bath, full electron-space ED tests the many-body algebra.
For a continuum local-moment problem, a converged NRG calculation can test
correlation physics that a few-level bath may miss. Match the Hamiltonian,
bandwidth, hybridization convention, temperature and observable before comparing
numbers. Use matched-bath comparisons to assess solver accuracy and independent
bath refinements to assess continuum accuracy.

\section{Canonical coordinates}\label{sec:coordinates}

All solvers act on the Hamiltonian written in \emph{canonical coordinates}: the impurity
operators $d_\sigma$ and the Bogoliubov quasiparticles $\gamma_{li\sigma}$ of the
\emph{isolated} reservoirs, \cref{eq:bogoliubov}. In these coordinates the reservoir
Hamiltonian is diagonal. When the QP expansion converges rapidly, the coupled
ground state is dominated by components with few bath quasiparticles: the
impurity is dressed by virtual quasiparticle excitations. Whether this description
is adequate must be checked at the couplings of interest. The
map from electron operators to canonical modes is stored with the model
(\code{metadata["coordinates"]}); \code{h.physical_operator} applies it to any electron
operator.

Two optional reductions exploit the symmetry of the single-dot junction. The eta
representation diagonalizes the particle--hole-type symmetry of \cref{sec:symmetries}.
Paired-mode compression keeps only the combination of mirror-symmetric levels that couples
to the dot and freezes the decoupled (\enquote{dark}) combinations in their vacuum. Both
depend on the parameters: the eta representation exists only at $\eps=0$, and the
compressed modes depend on $\phi$. They are therefore switched off for driven calculations
(\cref{sec:fixed-coordinates}).

\section{The quasiparticle expansion}\label{sec:qp}

The Fock space of the finite model is spanned by occupation-number states of all impurity
and quasiparticle modes. The QP expansion keeps every impurity configuration but at most
$q$ occupied quasiparticle modes,
\begin{equation}\label{eq:qp-projection}
  H_q=P_qHP_q ,\qquad
  P_q=\sum_{n\le q}\Pi_n ,
\end{equation}
where $\Pi_n$ projects onto states with exactly $n$ bath quasiparticles in the chosen
symmetry sector. Because $P_q$ projects onto a subspace, the lowest eigenvalue of $H_q$ is
a variational upper bound. It decreases monotonically with $q$ and reaches the exact
finite-model energy when $q$ equals the number of bath modes
(\code{cutoff=None})~\cite{ilicin2025variational,ilicin2026effects}. With $q=2$, the
singlet ground state is described by the dot configurations dressed by up to two
quasiparticles. This captures the Yu--Shiba--Rusinov binding and much of the Kondo
screening for moderate couplings. The dimension grows as $\binom{N}{q}$ with the number $N$
of bath modes; the code counts it exactly before building the basis. \Cref{tab:dimensions}
gives typical values for the reference junction.

The retained amplitudes are optimized together by diagonalization, summing
tunnelling processes to all orders within the selected QP space. The amplitudes
describe virtual QP components of a stationary state; nonequilibrium populations
and poisoning rates require a dynamical model. Each sector's ground energy is
an upper bound, while differences of independently truncated sector energies
can approach their converged values from either side. Check the excitation gaps
and transition locations themselves.

\begin{table}[htb]
\centering\small
\caption{Fock-space dimensions of the reference junction (two leads, \code{cosh_grid}),
even sector $S_z=0$ / odd sector $S_z=\tfrac12$.}\label{tab:dimensions}
\begin{tabular}{rrrrr}
\toprule
level pairs per lead & bath modes & $q=2$ & $q=3$ & $q=4$\\
\midrule
2 & 16 & 146 / 109 & 594 / 557 & 2\,162 / 1\,789 \\
4 & 32 & 546 / 409 & 4\,386 / 4\,249 & 33\,186 / 27\,609 \\
8 & 64 & 2\,114 / 1\,585 & 33\,858 / 33\,329 & 525\,890 / 438\,065 \\
16 & 128 & 8\,322 / 6\,241 & 266\,370 / 264\,289 & 8\,394\,882 / 6\,995\,041 \\
\bottomrule
\end{tabular}
\end{table}

A Floquet calculation multiplies these dimensions by the number of harmonic blocks, $2M+1$,
where $M$ is the Floquet cutoff (\cref{sec:sambe}). For large baths the static eigenbasis of
\cref{sec:sambe-method} is therefore useful.

\begin{incode}
\code{solve(h, cutoff, sector, options)} returns \code{Eigenstates} with \code{energies},
\code{vectors} (columns in the order of \code{basis}), \code{residuals}, the quasiparticle
number distribution \code{qp_weights} and the expectation values of all model
observables. The eigensolver (dense, sparse Lanczos or matrix-free) is chosen automatically
according to the dimension. \code{ProjectedOperator(op, basis)} applies an arbitrary operator
in the same Fock basis, also between different sectors.
The explicit \code{method="block"} uses locally optimal block preconditioned
conjugate gradients (LOBPCG), evolving several trial vectors together
to retain degenerate multiplets; compare the complete subspace with a dense or
symmetry-resolved reference. \code{QPWorkspace} reuses compatible preparation
without changing the finite Hamiltonian or its QP projection. For small problems,
its independent state count uses a full dense spectrum or Hermitian LDL factorization
to check how many states lie below a specified energy boundary.
\par\smallskip\usermanual{docs/interfaces.md}.
\end{incode}

\section{Bath representations}\label{sec:bath}

The finite bath is a quadrature of the continuous reservoir. The \code{cosh_grid}
discretization places levels with logarithmic density near the gap edge, where the bound
states draw their weight. Surrogate reservoirs fit a few levels to the hybridization
function on the imaginary axis~\cite{baran2023surrogate}, and the Pad\'e chain expansion
builds a short Wilson-like chain~\cite{bobok2025scalable}. In every case three distinct
approximations must be checked separately: the representation of the continuum by the
finite bath, the QP cutoff, and the eigensolver tolerance. A small residual measures only
the last of them.

\begin{caution}
The half-bandwidth $D$ is a physical parameter of the reservoir. Refining the bath at
fixed $D$ converges to a finite-band continuum. Approaching the wide-band limit
also requires increasing $D$ relative to the other physical energy scales.
\end{caution}

\section{Density-matrix renormalization group}\label{sec:dmrg}

The DMRG solver~\cite{white1992density,schollwock2011density} represents the many-body
state as a matrix-product state over the impurity and bath modes, and the Hamiltonian as a
matrix-product operator; both are built with the TeNPy library~\cite{hauschild2018efficient}.
DMRG retains all quasiparticle numbers and therefore complements the QP expansion when many
quasiparticles matter, for instance deep in the Kondo regime. Both solvers accept the same
\code{Hamiltonian} object, which allows direct comparisons.

An MPS controls the entanglement retained across partitions of the ordered orbital
list while including all quasiparticle numbers allowed by the sector. Its bond dimension controls
the number of states retained in a Schmidt decomposition. Increasing that dimension
allows more correlations, but the optimization must also converge and the relevant
excited states must be targeted. The orbital ordering affects the entanglement cost;
it specifies the tensor-product representation, while the Hamiltonian defines the physical geometry.
Schollw\"ock's review~\cite{schollwock2011density} provides the underlying
variational picture. A finite discretized bath remains an approximation with
either MPS or QP states.

\begin{physics}[Three convergence questions]
The eigenvector residual measures how accurately the chosen matrix problem was solved.
Increasing the QP cutoff or MPS bond dimension tests whether the interacting state
is represented adequately. Refining the bath tests whether that matrix problem
represents the intended reservoirs. Vary these controls separately and inspect
the observable of interest, especially near a parity crossing where small
relative errors can exchange the ground-state ordering.
\end{physics}

\section{Validation}

Calculations on small baths are compared with an independent exact diagonalization in the
physical electron basis, built from tensor products of single-mode operators without the
Bogoliubov transformation. Further checks use analytic limits, symmetry identities and
finite differences of energies, e.g.\ the current \eqref{eq:current} against the numerical
derivative $\partial E/\partial\phi$. The same independent diagonalization serves as the
reference for the driven problems (\cref{sec:validation-driven}).

\Cref{ch:practical} provides a complete runnable workflow using these methods,
including the static spectrum of \cref{fig:static} and the physical current.
